% STATUS: reviewed
\chapter{Laikui bėgant: veiksmų registras}\label{ch:register}

Atitikties ataskaita aprašo vieną dieną. Šis skyrius aprašo, ką sistema saugo tarp ataskaitų:
\term{veiksmų registrą} --- kiekvienos įmonės sąrašą veiksmų, kurių reikalavo jos patvirtintų
ataskaitų planai; kiekvienas veiksmas turi atsakingą asmenį, terminą, būseną ir išvadas, kurias
jis turi uždaryti. Skyrius išdėsto problemą ir vienintelį įvykį, kuris pajudina registrą
(\cref{sec:register:why,sec:register:event}), kas į registrą patenka per tą įvykį
(\cref{sec:register:plan}), veiksmo būsenas (\cref{sec:register:states}) ir --- skyriaus
šerdį --- kaip vėlesnės ataskaitos patvirtina ir uždaro veiksmus
(\cref{sec:register:settle}). Toliau aprašoma pareigos tapatybė, kuria visa tai remiasi,
išvados istorija ir teisės aktų pokyčiai
(\cref{sec:register:identity,sec:register:history,sec:register:law}); tada --- ką su registru
daro pokalbių asistentas ir Excel eksportas (\cref{sec:register:reading}), kur eina riba tarp
modelių, kodo ir žmonių (\cref{sec:register:line}) ir kas išbandyta su tikrais duomenimis
(\cref{sec:register:tried}).

Šiame skyriuje modelis niekur nekviečiamas. Todėl jame nėra kelio kortelės: registras yra kodas
ir žmonės, o \cref{sec:register:line} paaiškina, kodėl taip padaryta sąmoningai.


\section{Problema: planas, kuris baigdavo gyvenimą kartu su ataskaita}\label{sec:register:why}

Ataskaita baigiasi \term{veiksmų planu}: pagal prioritetą surikiuotu veiksmų sąrašu, kuriame
kiekvienas veiksmas nurodo išvadas, kurias turi uždaryti, laikotarpį ir atsakingo asmens rolę
(kaip planas sudaromas, aprašyta \cref{ch:report}). Planas saugomas kartu su ataskaita, lentelėje
\code{report_actions}. Kita tos pačios įmonės analizė parengia naują ataskaitą su nauju planu, ir
niekas jų nesieja: niekas negali pasakyti, ar senojo plano veiksmas buvo įvykdytas, kas jį
įvykdė, iki kada ir ar spraga, kurią jis šalino, išnyko.

Kurso išvadų ir veiksmų registras į tai atsako viena taisykle: veiksmas uždaromas tik po
pakartotinės patikros. Nepakanka pasakyti, kad darbas atliktas; kažkas turi pažiūrėti dar kartą
ir rasti reikalavimą įvykdytą. Sistema jau turi būdą pažiūrėti dar kartą --- kitą ataskaitą ---
ir žmogų, kuris prisiima atsakomybę už ataskaitą, --- tą, kuris ją patvirtina. Registras šiuos
du dalykus sujungia:

\begin{itemize}
  \item kiekvienos patvirtintos ataskaitos planą jis saugo kaip savo eilutes, atskirai kiekvienai
    įmonei, lentelėje \code{company_actions} (modelis \code{CompanyAction}), ir šios eilutės
    išlieka ilgiau nei ataskaita;
  \item kitą patvirtintą ataskaitą jis naudoja kaip pakartotinę patikrą: veiksmas uždaromas, kai
    vėlesnės patvirtintos ataskaitos randa įvykdytas visas jo apimamas pareigas.
\end{itemize}

Visa tai yra viename API faile, \file{analysis/register.py}, su septyniomis viešomis funkcijomis:
\code{sign_off}, \code{withdraw}, \code{standing}, \code{state_of}, \code{history},
\code{changes} ir \code{law_changes}. Penkios konstantos failo pradžioje nustato viso skyriaus
žodyną.

\begin{table}[H]
\caption{Konstantos registro failo pradžioje}\label{tab:register:constants}
\begin{tabularx}{\textwidth}{@{}L{22mm}L{58mm}Y@{}}
\toprule
\hd{Konstanta} & \hd{Reikšmė} & \hd{Kam naudojama} \\
\midrule
\code{DUE_DAYS} & \code{immediate} 14, \code{30_days} 30, \code{90_days} 90, \code{12_months} 365 &
  Dienų skaičius nuo patvirtinimo iki veiksmo termino pagal plano laikotarpį. \\
\code{OPEN} & \code{missing}, \code{partial}, \code{unknown} & Išvados būsenos, kurioms reikia veiksmo. \\
\code{SETTLED} & \code{met}, \code{not_applicable} & Būsenos, kurios išvadą patvirtina. \\
\code{LIVE} & \code{open}, \code{in_progress}, \code{done} & Neuždaryto veiksmo būsenos. \\
\code{HISTORY} & 4 & Kiek ankstesnių vertinimų siekia išvados istorija. \\
\bottomrule
\end{tabularx}
\end{table}

\code{OPEN} ir \code{SETTLED} penkias išvados būsenas suskirsto į dvi grupes, kuriomis registras
samprotauja. Išvada, kurios būsena \status{missing}, \status{partial} ar \status{unknown}, yra
\emph{atvira}: jai reikia veiksmo, o vėlesnė ataskaita, kuri ją randa tokią pat, nieko
nepatvirtina. Išvada, kurios būsena \status{met} ar \status{na}, yra \emph{išspręsta}. Kur tik
registras skaito būseną, jis skaito \emph{galiojančią} būseną --- peržiūrėtojo sprendimą, jei
toks yra, kitaip modelio vertinimą (\code{_effective}). \code{LIVE} įvardija tris neuždaryto
veiksmo būsenas.


\section{Vienintelis įvykis, kuris pajudina registrą}\label{sec:register:event}

Registras keičia savo patvirtinimus tik per vieną įvykį: kai žmogus \term{patvirtina baigtą
ataskaitą} arba kai tas patvirtinimas atšaukiamas. Patvirtinimas --- tai maršrutas
\route{PATCH}{/api/reports/{id}} su turiniu \code{{approved: true, note?}}; konsolė jį siunčia
mygtuku \ui{Patvirtinti ataskaitą} ataskaitos puslapyje, o jo dialogas dar prieš patį mygtuką
pasako, ką patvirtinimas padarys registrui
(\file{pages/Reports/ReportApproval.jsx}). Patvirtinti galima tik baigtą ataskaitą (kitaip ---
409 \code{NOT_FINISHED}). Rodinys \code{_approve} faile \file{analysis/api/views.py} ataskaitoje
įrašo \code{approved_by}, \code{approved_at} ir \code{approval_note}, o tada iškviečia registrą;
komentaras toje eilutėje sako: \enquote{Įmonės registras pajuda patvirtinus ataskaitą ir dėl nieko
kito}.

\begin{figure}[H]
\centering
\begin{tikzpicture}[node distance=3.5mm and 10mm]
  \node[man, text width=41mm] (sign) {patvirtina baigtą ataskaitą};
  \node[man, text width=41mm, below=of sign] (unsign) {atšaukia patvirtinimą};
  \node[man, text width=41mm, below=of unsign] (review) {pakeičia patvirtintos ataskaitos išvadą};
  \node[man, text width=41mm, below=of review] (delete) {ištrina patvirtintą ataskaitą};
  \coordinate (mid) at ($(unsign.east)!.5!(review.east)$);
  \node[chk, text width=37mm, right=10mm of mid] (settle) {\code{_settle}\\ visi išvadų patvirtinimai iš naujo apskaičiuojami pagal visas patvirtintas ataskaitas};
  \node[chk, text width=30mm, right=of settle, yshift=13mm] (plan) {\code{sign_off}, 2~žingsnis\\ tik patvirtinant:\\ įtraukiamas ataskaitos planas};
  \node[store, inner xsep=2pt, below=9mm of plan] (table) {\code{company_actions}};
  \node[man, text width=41mm, below=7mm of delete] (edit) {keičia vieną veiksmą: atsakingą asmenį, terminą, būseną, pastabą};
  \draw[flow] (sign.east) -- ++(4mm,0) |- ([yshift=3mm]settle.west);
  \draw[flow] (unsign.east) -- ++(4mm,0) |- ([yshift=1mm]settle.west);
  \draw[flow] (review.east) -- ++(4mm,0) |- ([yshift=-1mm]settle.west);
  \draw[flow] (delete.east) -- ++(4mm,0) |- ([yshift=-3mm]settle.west);
  \draw[flow] (settle.east) -- ++(4mm,0) |- (plan.west);
  \draw[flow] (settle.east) -- ++(4mm,0) |- (table.west);
  \draw[flow] (plan.south) -- (table.north);
  \draw[flow] (edit.east) -| (table.south) node[lbl, pos=.25] {\code{PATCH}~~\code{/api/actions/{id}}};
\end{tikzpicture}
\actorlegend
\caption{Viskas, kas rašo į registrą. Keturi dalykai, kuriuos žmogus daro su ataskaita, pasiekia vienintelę
funkciją, kuri apskaičiuoja patvirtinimus; tik patvirtinimas dar ir įtraukia planą. Žmogus gali
tiesiogiai keisti atskirą veiksmą. Nei užduotis, nei rengiama ar baigiama ataskaita, nei
įkeltas dokumentas neturi kelio į šią lentelę, o pokalbis neturi įrankio, kuris į ją rašytų:
peržiūra, kurią jis įrašo naudotojo nurodymu, yra trečiasis iš čia parodytų dalykų
(\cref{sec:register:reading}).}\label{fig:register:events}
\end{figure}

\Cref{fig:register:events} rodo visus, kas rašo į lentelę. Dvi \file{analysis/register.py}
funkcijos kviečiamos iš rodinių ir niekur kitur kode:

\begin{table}[H]
\caption{Įvykiai, kurie pajudina registrą, ir ką kiekvienas iškviečia}\label{tab:register:events}
\begin{tabularx}{\textwidth}{@{}L{48mm}L{37mm}Y@{}}
\toprule
\hd{Įvykis} & \hd{Iškviečia} & \hd{Poveikis} \\
\midrule
Žmogus patvirtina baigtą ataskaitą (\code{_approve}, \code{approved: true}) &
  \code{register.sign_off} &
  Esami veiksmai patikrinami pagal visas patvirtintas ataskaitas, dabar ir pagal šią; tada
  įtraukiamas ataskaitos planas. Grąžina \code{{added, covered, confirmed, closed, reopened}}. \\
Žmogus atšaukia patvirtinimą (\code{_approve}, \code{approved: false}) &
  \code{register.withdraw} &
  Patvirtinimai iš naujo apskaičiuojami be šios ataskaitos. Grąžina
  \code{{withdrawn}} --- pasikeitusių veiksmų skaičių. \\
Peržiūrėtojas pakeičia patvirtintos ataskaitos išvadą (\code{finding_item}) &
  \code{register.withdraw} &
  Bet koks išvados pakeitimas --- pakeista būsena, peržiūrėtojo pastaba, priežastis ---
  panaikina patvirtinimą: pakeistam vertinimui reikia naujo patvirtinimo. Kartu su juo atšaukiama ir
  tai, ką jis buvo patvirtinęs. \\
Ištrinama patvirtinta ataskaita (\code{report_item}, \code{DELETE}) &
  \code{register.withdraw} &
  Pirmiausia panaikinamas patvirtinimas, registras apskaičiuojamas be ataskaitos, tada ataskaita
  ištrinama: ji nusineša tai, ką patvirtino tik ji viena. \\
\bottomrule
\end{tabularx}
\end{table}

Atsakymas į patvirtinimą grąžina ataskaitą ir, po raktu \code{register}, penkis skaičius; konsolė
juos pateikia sėkmės pranešime su nuoroda į registrą arba pasako, kad registras nepasikeitė, kai
visi penki lygūs nuliui. Ir patvirtinimas, ir jo atšaukimas įrašomi į veiksmų žurnalą
(\code{reports.approve}, \code{reports.unapprove}) kartu su tais skaičiais.

\begin{humanbox}{Kodėl jį pajudina tik parašas}
Baigta ataskaita yra modelio įvertinimas, kurio žmogus dar nepriėmė. Jos išvadas dar galima
pakeisti peržiūroje, o modelis \code{Report} patvirtinimą aprašo kaip žmogaus pritarimą tam, kad
\enquote{peržiūrėta ataskaita yra įmonės įvertinimas}. Jei veiksmą galėtų uždaryti nepatvirtinta
ataskaita, modelio vertinimas nutrauktų žmogaus prisiimtą įsipareigojimą. Todėl įrodymas, kuris
uždaro veiksmą, visada yra ataskaita, po kuria kažkas pasirašė, o kai tas parašas atšaukiamas ---
ranka arba todėl, kad vėliau buvo pakeista išvada, --- kartu nyksta ir uždarymas. Pačio registro
failo antraštė tai pasako trumpai: uždarymas yra pakartotinė patikra, o ne pareiškimas.
\end{humanbox}

Tas pats samprotavimas atmeta ir kitus kandidatus. Nauja, dar nepatvirtinta ataskaita registre
niekur neatsiranda; testas \code{test_a_sign_off_brings_the_plan_in_once} tikrina, kad
nepatvirtinta ataskaita \enquote{nieko nepajudina}. Nė viena užduoties apdorojimo funkcija
neimportuoja registro rašymo funkcijų. Pokalbių asistentas registrui turi vieną įrankį, ir tas
tik skaito (\cref{sec:register:reading}). Įkėlus dokumentą, kuris uždarytų spragą, niekas
nepasikeičia, kol vėlesnė ataskaita jo neįvertina ir žmogus tos ataskaitos nepatvirtina.


\section{Kas patenka į registrą patvirtinus ataskaitą}\label{sec:register:plan}

\code{sign_off(report, user)} vykdoma vienoje duomenų bazės transakcijoje, dviem žingsniais.
1~žingsnis esamus veiksmus patikrina pagal patvirtintas ataskaitas (\cref{sec:register:settle});
2~žingsnis, aprašomas čia, įtraukia pačios ataskaitos planą. Jis eina per plano veiksmus jų
eilės tvarka ir kiekvienam sudaro \term{elementų} sąrašą: išvadų, kurias veiksmas turi uždaryti.

Plano veiksmo išvada tampa elementu, nebent tenkinama viena iš trijų sąlygų:

\begin{enumerate}
  \item \textbf{Ji nėra atvira.} Jos galiojanti būsena nepatenka į \code{OPEN}: tam, kas
    įvykdyta, veiksmo nereikia.
  \item \textbf{Naujesnė patvirtinta ataskaita ją jau randa išspręstą.} \code{_later(report)}
    iš įmonės patvirtintų ataskaitų, parengtų \emph{po} šios, surenka naujausią kiekvienos
    pareigos vertinimą; kai tas vertinimas patenka į \code{SETTLED}, nebelieka ko daryti. Tai
    svarbu, kai senesnė ataskaita patvirtinama pavėluotai.
  \item \textbf{Jos pareigą jau apima neuždarytas veiksmas.} Aibėje \code{under_way} yra kiekviena
    pareiga, kurią nurodo bet kurio \code{LIVE} būsenos veiksmo elementas. Tokia išvada
    skaičiuojama kaip \code{covered} ir praleidžiama: pareigos niekada neapima du
    neuždaryti veiksmai. Šiuo patvirtinimu įtraukta pareiga į aibę patenka iš karto, todėl jos
    negali perimti abu to paties plano veiksmai.
\end{enumerate}

Elementas yra momentinė kopija, saugoma veiksmo JSON stulpelyje \code{items}:
\code{obligationId}, \code{findingId}, \code{title} (pareigos pavadinimas, kaip jį pateikia
išvada), \code{act} ir \code{path} (išvados pirmosios nuorodos aktas, o jei jos nėra --- pačios
išvados aktas, ir tos nuorodos nuostata), \code{sourceKind}, \code{status} (galiojanti būsena
patvirtinimo metu) ir \code{confirmed}, kuris iš pradžių yra null. Kaip ir ataskaita, registras
nusikopijuoja rodomus pavadinimus, todėl veiksmas atrodo taip pat ir pasikeitus teisės aktų
bazei ar ataskaitai.

Plano veiksmas, kurio visos išvados buvo praleistos, neįtraukiamas. Kitu atveju sukuriamas
\code{CompanyAction} su plano \code{title}, \code{description}, \code{category},
\code{priority} ir \code{effort}; plano atsakingo asmens role kaip \code{owner}; būsena
\code{open}; ir terminu, skaičiuojamu nuo patvirtinimo momento pagal plano laikotarpį
(\code{DUE_DAYS}): 14 dienų, kai laikotarpis yra \code{immediate}, 30, kai \code{30_days}, 90,
kai \code{90_days}, 365, kai \code{12_months}, ir 90, kai laikotarpio lentelė neįvardija. Terminas
skaičiuojamas nuo patvirtinimo, o ne nuo ataskaitos parengimo dienos, nes patvirtinimo diena yra
ta diena, kai žmogus priėmė planą.

\paragraph{Nieko neįtraukiama du kartus.} Kiekvienas registro veiksmas prisimena plano veiksmą,
iš kurio kilo (\code{plan_action_id}), o įmonės ir šio identifikatoriaus unikalumo
apribojimas paremia patikrą kode. Todėl atšauktas ir vėl suteiktas patvirtinimas nieko neprideda; neprideda
ir kita ataskaita, kurios planas apima tas pačias atviras pareigas, --- jos išvados
skaičiuojamos kaip \code{covered}.

\paragraph{Veiksmai, kurie neuždaro jokios išvados.} Plane gali būti veiksmų be išvadų ---
ištaisyti nuorodas, kurias teisės aktų nuorodų patikra rado klaidingas, išsiaiškinti faktą,
kurio įmonės profilis nenurodo. Jie taip pat įtraukiami, su tuščiu sąrašu \code{items}. Neturėdami
pareigos, pagal kurią būtų atpažįstami, jie sutapatinami pagal savo \code{category}: kol
\emph{kitos} ataskaitos veiksmas be elementų toje kategorijoje dar neuždarytas, vėlesnės ataskaitos
veiksmas be išvadų toje pačioje kategorijoje yra tas pats dar kartą paprašytas darbas, jis
skaičiuojamas kaip \code{covered} ir neįtraukiamas. Jokia ataskaita tokio veiksmo niekada negali
patvirtinti; jį uždaro tik žmogus (\cref{sec:register:states}).


\section{Veiksmo būsenos ir ką gali pakeisti žmogus}\label{sec:register:states}

Veiksmas saugomas vienoje iš keturių būsenų, \code{ACTION_STATES} faile
\file{analysis/models.py}, o jų laikymąsi užtikrina duomenų bazės tikrinimo apribojimas.

\begin{table}[H]
\caption{Veiksmo būsenos. Pirmosios keturios saugomos; penktoji apskaičiuojama skaitant registrą}\label{tab:register:states}
\begin{tabularx}{\textwidth}{@{}L{22mm}L{45mm}Y@{}}
\toprule
\hd{Būsena} & \hd{Konsolėje} & \hd{Reiškia} \\
\midrule
\code{open} & \ui{Atviras} & Įtrauktas patvirtinus ataskaitą; apie jį dar nieko nepranešta. \\
\code{in_progress} & \ui{Vykdomas} & Kažkas prie jo dirba. \\
\code{done} & \ui{Įvykdytas, laukia patikros} &
  Atsakingo asmens žodis, kad veiksmas atliktas. Tai ne uždarymas: veiksmas laukia vėlesnės
  patvirtintos ataskaitos. Nustačius šią būseną, įrašomas \code{done_at}. \\
\code{closed} & \ui{Uždarytas} &
  Arba vėlesnės patvirtintos ataskaitos rado įvykdytas visas jo apimamas pareigas
  (\code{confirmed_report} nurodo uždariusią ataskaitą), arba žmogus jį uždarė ranka, nurodęs
  priežastį. \\
\midrule
\code{late} & \ui{Vėluojantis} &
  Niekada nesaugoma. \code{state_of(action, today)} atsako \code{late} veiksmui, kuris yra
  \code{open} ar \code{in_progress} ir kurio termino data ankstesnė už šiandienos. \\
\bottomrule
\end{tabularx}
\end{table}

Vėlavimas apskaičiuojamas, nes jis priklauso nuo kalendoriaus, o ne nuo to, kas nutiko
veiksmui: saugomas požymis jau kitą rytą po termino būtų klaidingas, nebent jį perrašytų kokia
nors užduotis. Veiksmas, pažymėtas \code{done}, niekada nerodomas kaip vėluojantis: atsakingas
asmuo atsakė, ir veiksmas laukia pakartotinės patikros. Kiekvienas API atsakymas grąžina abi
reikšmes: \code{status} --- saugomą būseną ir \code{state} --- rodytiną būseną.

\begin{figure}[H]
\centering
\begin{tikzpicture}[x=1mm, y=1mm]
  \node[man, text width=20mm] (in) at (11,0) {ataskaita patvirtinama};
  \node[svc, minimum width=20mm] (open) at (47,0) {\code{open}};
  \node[svc, minimum width=26mm] (prog) at (84,0) {\code{in_progress}};
  \node[svc, minimum width=20mm] (done) at (121,0) {\code{done}};
  \begin{scope}[on background layer]
    \node[lane, fit=(open)(prog)(done), inner sep=7pt] (live) {};
  \end{scope}
  \node[draw=warn, dashed, rounded corners=3pt, line width=.6pt, fit=(open)(prog), inner sep=3.5pt] (late) {};
  \node[lanelabel, anchor=south west] at (live.north west) {neuždarytos būsenos};
  \node[note, text=warn, anchor=north, text width=24mm, align=center] at (11,-8)
    {brūkšninis rėmelis: praėjus terminui rodoma \code{late}};
  \draw[flow] (in) -- (open) node[lbl, pos=.4, above=1pt] {planas};
  \draw[flow, <->] (open) -- (prog) node[lbl, midway, above=1.5pt] {žmogus};
  \draw[flow, <->] (prog) -- (done) node[lbl, midway, above=1.5pt] {žmogus};
  \draw[flow] (done.north) -- ++(0,9) -| (prog.north)
    node[lbl, pos=.25, above=1pt] {patvirtinimas ar jo atšaukimas: po \code{done_at} parengta ataskaita vis dar randa atvirą išvadą};
  \node[chk, text width=38mm] (crep) at (58,-31) {\code{closed}: uždarė ataskaita\\ nustatytas \code{confirmed_report}};
  \node[man, text width=38mm] (chand) at (116,-31) {\code{closed}: uždarė žmogus\\ pastaba nurodo kodėl};
  \coordinate (r1) at (52,0);  \coordinate (r2) at (64,0);
  \coordinate (h1) at (110,0); \coordinate (h2) at (122,0);
  \draw[flow] (r1 |- live.south) -- (r1 |- crep.north)
    node[lbl, midway, left=2pt, align=right] {patvirtinimas: visos\\ išvados patvirtintos};
  \draw[opt] (r2 |- crep.north) -- (r2 |- live.south)
    node[lbl, midway, right=2pt, align=left] {įrodymas nebegalioja:\\ atgal į \code{closed_from}};
  \draw[flow] (h1 |- live.south) -- (h1 |- chand.north)
    node[lbl, midway, left=2pt, align=right] {žmogus,\\ su pastaba};
  \draw[flow] (h2 |- chand.north) -- (h2 |- live.south)
    node[lbl, midway, right=2pt, align=left] {žmogus\\ atšaukia};
\end{tikzpicture}
\actorlegend
\caption{Veiksmo būsenos ir kas veiksmą perkelia iš vienos į kitą. Žmogus laisvai perkelia veiksmą
tarp neuždarytų būsenų ir gali jį uždaryti ranka; kodas veiksmą uždaro, grąžina į darbą ar vėl
atidaro tik patvirtinus ataskaitą ar atšaukus patvirtinimą. Jokia rodyklė iš ataskaitos atlikto
uždarymo neveda ranka.}\label{fig:register:states}
\end{figure}

\subsection{Žmogus keičia vieną veiksmą}

Tarp patvirtinimų žmogus atnaujina veiksmo duomenis per maršrutą
\route{PATCH}{/api/actions/{id}}
(\code{action_item}); užklausą siunčia atvertos eilutės forma įmonės skiltyje \ui{Veiksmai}.
Turinyje yra tik tai, kas pasikeitė.

\begin{table}[H]
\caption{Ką žmogus gali pakeisti veiksme}\label{tab:register:patch}
\begin{tabularx}{\textwidth}{@{}L{20mm}Y@{}}
\toprule
\hd{Laukas} & \hd{Taisyklė} \\
\midrule
\code{owner} & Tekstas be pradinių ir galinių tarpų, ne daugiau kaip 300 ženklų. Iš pradžių tai
  plano atsakingo asmens rolė; žmogus ją pakeičia tuo, kas iš tikrųjų atsakingas už darbą. \\
\code{dueDate} & Data, \code{YYYY-MM-DD}. \\
\code{note} & Tekstas be pradinių ir galinių tarpų arba null, ne daugiau kaip 4\,000 ženklų. \\
\code{status} & Viena iš \code{open}, \code{in_progress}, \code{done}, \code{closed} --- niekada
  \code{late}. Pakeitus į \code{done}, įrašomas \code{done_at}. Uždarant reikia pastabos. Vėl
  atidaryti galima tik ranka uždarytą veiksmą. \\
\bottomrule
\end{tabularx}
\end{table}

Klaidinga reikšmė atmetama su 400 \code{INVALID_FIELD}; naudotojas, kuris nemato veiksmo įmonės,
gauna 404. Kiekvienas pakeitimas įrašomas į veiksmų žurnalą kaip \code{actions.update} su
pasikeitusiais laukais (pastabos tekstas ten nekopijuojamas, tik pats faktas, kad ji pasikeitė).
Nėra maršruto, kuris sukurtų veiksmą, ir nėra tokio, kuris jį ištrintų: veiksmai atsiranda iš planų ir
išnyksta tik kartu su savo įmone.

\begin{humanbox}{Uždarymas ranka ir ko ranka negali atšaukti}
Žmogus gali uždaryti veiksmą be jokios ataskaitos patvirtinimo --- paslauga nutraukta, pareiga
įmonei nebeaktuali, --- bet tik su pastaba: be jos API atsako 400 \code{NOTE_REQUIRED}
(\enquote{Uždarant veiksmą be ataskaitos patvirtinimo, pastaboje nurodykite priežastį}), o
konsolė laiko išsaugojimo mygtuką išjungtą, kol pastaba neparašyta. Tada veiksme įrašoma, kas ir
kada jį uždarė (\code{closed_by}, \code{closed_at}), ir būsena, iš kurios jis buvo uždarytas, o
registras uždarymą rodo kaip \code{manual}. Žmogus gali jį atšaukti, nustatydamas kitą būseną;
tada uždarymo laukai išvalomi.

\smallskip
Ataskaitos atliktas uždarymas yra kas kita. Bandymas nustatyti jam kitą būseną atmetamas su 409
\code{CLOSED_BY_REPORT}: \enquote{Šį veiksmą uždarė patvirtinta ataskaita --- norėdami jį
ginčyti, peržiūrėkite išvadą toje ataskaitoje}. Uždarymas remiasi patvirtintos ataskaitos
išvadomis, todėl ginčyti jį reikia ties išvada. Ten pakeitus išvadą, ataskaitos patvirtinimas
atšaukiamas, o tas pats kodas, kuris veiksmą uždarė, apskaičiuoja uždarymą iš naujo be tos
ataskaitos (\cref{sec:register:settle}). Tokio veiksmo būsenos pasirinkimo lauką konsolė rodo išjungtą,
o po juo --- eilutę, kuri pasako, kur uždarymą ginčyti.
\end{humanbox}

\subsection{Registro skaitymas}

\route{GET}{/api/companies/{id}/actions} (\code{company_actions}) grąžina registrą ir veiksmų
skaičių kiekvienoje rodomoje būsenoje. Tvarka yra tokia, kokios reikia žmogui: pirma vėluojantys
veiksmai, tada atviri ir vykdomi, tada pažymėti kaip įvykdyti, tada uždaryti; grupės viduje ---
pagal terminą, tada pagal plano prioritetą. Sąrašas apribojamas 500 eilučių; skaičiai ---
ne. Skiltis, kuri jį rodo, aprašyta \cref{ch:console}.


\section{Patvirtinimas vėlesnėmis ataskaitomis}\label{sec:register:settle}

Šis poskyris yra registro šerdis: kaip vėlesnė ataskaita patvirtina išvadą, uždaro veiksmą ar
grąžina jį į darbą. Visa tai atlieka viena privati funkcija, \code{_settle(company_id,
now)}, ir ją kviečia tiek \code{sign_off}, tiek \code{withdraw}. Ji netaiko naujai patvirtintos
ataskaitos registrui kaip pakeitimo. Kiekvieną kartą ji visus patvirtinimus apskaičiuoja iš
naujo iš \emph{visų} patvirtintų įmonės ataskaitų. Iš šios savybės išplaukia visa kita: esamų
veiksmų patvirtinimai ir uždarymai išeina tokie patys, kad ir kokia tvarka buvo suteikti parašai,
o parašui atšaukti nereikia jokio atšaukimų žurnalo --- funkcija tiesiog vykdoma dar kartą be tos ataskaitos.

\subsection{Ką funkcija skaito}

\code{_settle} ima tuos įmonės veiksmus, kurie turi elementų ir yra arba neuždaryti, arba uždaryti
ataskaitos. Žmogaus ranka uždaryti veiksmai ir veiksmai be elementų paliekami ramybėje. Elementų
nurodytoms pareigoms ji įkelia visas jų išvadas iš įmonės ataskaitų, kurios yra baigtos ir
patvirtintos, surikiuotas pagal tai, kada ataskaita buvo \emph{parengta} (\code{created_at}), ir
kiekvienai pareigai išsaugo vertinimų sąrašą: kada ataskaita parengta, kuri tai ataskaita ir
galiojanti būsena.

Vienam veiksmui skaičiuojami tik tų ataskaitų vertinimai, kurios parengtos ne anksčiau nei paties
veiksmo ataskaita (\code{since}; jei ta ataskaita ištrinta --- ne anksčiau nei pats veiksmas). Ataskaita, parengta dar prieš atsirandant veiksmui, nieko nepasako apie
pagal jį atliktą darbą, net jei ji patvirtinama vėliau.

\subsection{Taisyklės}

Kiekvienam veiksmo elementui funkcija eina per pareigos vertinimus ta tvarka, kuria buvo rengiamos
ataskaitos, ir ieško ataskaitos, kuri \emph{pirmoji rado pareigą išspręstą ir po kurios nė viena
vėlesnė ataskaita jos nerado atviros}: išspręstas vertinimas pradeda seką, atviras vertinimas ją
nutraukia.

\begin{enumerate}
  \item \textbf{Išvada patvirtinama}, kai naujausias jos pareigos vertinimas yra išspręstas.
    Elemento laukui \code{confirmed} priskiriama \code{{reportId, at}}, nurodanti pirmąją tos
    nenutrūkusios sekos ataskaitą. Ataskaita, kuri pareigos nevertino, vertinimo nepalieka, todėl
    nieko nepatvirtina ir nieko nenutraukia: tarpinė siaura ataskaita apie kelis dokumentus
    registro netrikdo. \status{unknown} vertinimas yra atviras --- ataskaita, kuri negalėjo rasti
    įrodymų, nepatvirtina.
  \item \textbf{Veiksmas uždaromas}, kai patvirtintas kiekvienas jo elementas. Jo būsena tampa
    \code{closed}, \code{closed_from} prisimena ankstesnę būseną, \code{closed_by} lieka tuščias,
    o \code{confirmed_report} nurodo naujausią iš patvirtinusių ataskaitų --- tą, kuri rinkinį
    užbaigė.
  \item \textbf{Įvykdytu pažymėtas veiksmas grąžinamas į darbą} (\code{in_progress}), kai
    naujausias vienos iš jo pareigų vertinimas yra atviras \emph{ir} kilęs iš ataskaitos,
    parengtos po to momento, kai atsakingas asmuo pažymėjo veiksmą įvykdytu (\code{done_at}).
    Atsakingas asmuo pasakė, kad darbas atliktas; ataskaita, kuri pažiūrėjo vėliau, spragą vis
    tiek rado.
  \item \textbf{Įvykdytu pažymėto veiksmo negrąžina} ataskaita, parengta prieš \code{done_at}, net
    jei ta ataskaita patvirtinama vėliau: ji niekada nematė baigto darbo. Jo negrąžina ir
    ataskaita, kuri pareigos nevertino. Veiksmo, kuris yra \code{open} ar \code{in_progress},
    atviras vertinimas niekada nepajudina; jis tiesiog lieka, koks buvo.
  \item \textbf{Uždarytas veiksmas lieka uždarytas.} Ataskaitos uždarytam veiksmui skaitomi tik
    vertinimai iki uždariusios ataskaitos, todėl vėlesnė ataskaita, vėl radusi pareigą atvirą, jo
    neatidaro. Pačios tos ataskaitos planas spragai atneša \emph{naują} veiksmą, nes pareigos
    nebeapima joks neuždarytas veiksmas (\cref{sec:register:plan}), o kai ji randa pareigą
    \status{missing} ar \status{partial}, ataskaita tą spragą skaičiuoja kaip vėl atsiradusią
    (\cref{sec:register:history}). Senasis veiksmas lieka teisingu įrašu: spraga tada buvo
    uždaryta, o vėliau vėl atsivėrė.
  \item \textbf{Uždarymas nebegalioja tik dingus jį pagrindžiančiam įrodymui.} Kol uždariusi ataskaita
    patvirtinta, elementai skaitomi pagal ataskaitas iki jos; kai nebe --- pagal kiekvieną
    ataskaitą, kuri vis dar patvirtinta. Jei kiekvienas elementas tebėra patvirtintas, veiksmas
    lieka uždarytas, o \code{confirmed_report} nurodo naujausią likusią patvirtinusią ataskaitą.
    Jei bent vienas nebepatvirtintas, uždarymas nebegalioja: veiksmas grįžta į būseną, iš kurios buvo
    uždarytas (\code{closed_from}), \code{confirmed_report} ir \code{closed_at} išvalomi, o
    elementai dar kartą perskaitomi pagal kiekvieną patvirtintą ataskaitą.
\end{enumerate}

Funkcija grąžina keturis skaičius --- naujai patvirtintas išvadas, uždarytus veiksmus, į darbą
grąžintus veiksmus ir apskritai paliestus veiksmus. Patvirtinimas pirmuosius tris praneša kaip
\code{confirmed}, \code{closed} ir \code{reopened}; atšaukimas ketvirtąjį praneša kaip
\code{withdrawn}.

\subsection{Pavyzdys laiko juostoje}

\Cref{fig:register:timeline} seka vieną veiksmą per tris ataskaitas. Tai pavyzdys, šiam skyriui
sudarytas pagal pirmiau išdėstytas taisykles, o ne tikrų duomenų įrašas; testas
\code{test_a_later_signed_report_is_the_re_test} faile \file{tests/test_register.py} tokią pačią
seką atlieka per API.

\begin{figure}[H]
\centering
\begin{tikzpicture}[x=1mm, y=1mm,
    cell/.style={align=center, text width=23mm, inner sep=1.5pt, font=\sffamily\scriptsize, text=body},
    rowlabel/.style={align=right, text width=19mm, anchor=east, inner sep=1.5pt,
                     font=\sffamily\scriptsize\bfseries, text=ink},
    ev/.style={man, text width=20.5mm, inner xsep=3pt, inner ysep=3pt, minimum height=11mm, font=\sffamily\scriptsize},
    st/.style={chk, text width=20.5mm, inner xsep=3pt, inner ysep=3pt, minimum height=7mm, font=\sffamily\scriptsize}]
  % columns
  \foreach \x/\n in {33.5/1, 58.5/2, 83.5/3, 108.5/4, 133.5/5} {\coordinate (c\n) at (\x,0);}
  % time arrow
  \draw[flow] (21,0) -- (147,0);
  \node[rowlabel] at (20,0) {Šia tvarka};
  \foreach \n in {1,3,4} {\fill[ink] (c\n) circle[radius=.9];}
  \node[note, anchor=south] at (33.5,1) {parengta 1~ataskaita};
  \node[note, anchor=south] at (83.5,1) {parengta 2~ataskaita};
  \node[note, anchor=south] at (108.5,1) {parengta 3~ataskaita};
  % events
  \node[rowlabel] at (20,-10) {Žmogus};
  \node[ev] at (33.5,-10) {patvirtina 1~ataskaitą};
  \node[ev] at (58.5,-10) {pažymi veiksmą įvykdytu};
  \node[ev] at (83.5,-10) {patvirtina 2~ataskaitą};
  \node[ev] at (108.5,-10) {patvirtina 3~ataskaitą};
  \node[ev] at (133.5,-10) {atšaukia 3~ataskaitos patvirtinimą};
  % obligation X
  \node[rowlabel] at (20,-24) {Pareiga X};
  \node[cell] at (33.5,-24) {\status{missing}\\ veiksmo elementas};
  \node[cell] at (58.5,-24) {---};
  \node[cell] at (83.5,-24) {\status{met}\\ patvirtino 2~ataskaita};
  \node[cell] at (108.5,-24) {\status{met}\\ tebegalioja 2~ataskaitos patvirtinimas};
  \node[cell] at (133.5,-24) {patvirtino 2~ataskaita};
  % obligation Y
  \node[rowlabel] at (20,-38) {Pareiga Y};
  \node[cell] at (33.5,-38) {\status{partial}\\ veiksmo elementas};
  \node[cell] at (58.5,-38) {---};
  \node[cell] at (83.5,-38) {\status{partial}\\ nepatvirtinta};
  \node[cell] at (108.5,-38) {\status{met}\\ patvirtino 3~ataskaita};
  \node[cell] at (133.5,-38) {patvirtinimas panaikintas};
  % state
  \node[rowlabel] at (20,-52) {Veiksmas};
  \node[st] (s1) at (33.5,-52) {\code{open}};
  \node[ev, minimum height=7mm] (s2) at (58.5,-52) {\code{done}};
  \node[st] (s3) at (83.5,-52) {\code{in_progress}};
  \node[st] (s4) at (108.5,-52) {\code{closed}};
  \node[st] (s5) at (133.5,-52) {\code{in_progress}};
  \draw[flow] (s1) -- (s2);  \draw[flow] (s2) -- (s3);  \draw[flow] (s3) -- (s4);  \draw[flow] (s4) -- (s5);
  \node[cell, text=muted] at (33.5,-61) {įtraukiamas planas; terminas skaičiuojamas nuo čia};
  \node[cell, text=muted] at (58.5,-61) {įrašomas \code{done_at}};
  \node[cell, text=muted] at (83.5,-61) {grąžintas: Y atvira po \code{done_at} parengtoje ataskaitoje};
  \node[cell, text=muted] at (108.5,-61) {abu elementai patvirtinti; uždarė 3~ataskaita};
  \node[cell, text=muted] at (133.5,-61) {grįžta į būseną, iš kurios buvo uždarytas};
  % separators
  \begin{scope}[on background layer]
    \foreach \x in {46, 71, 96, 121} {\draw[draw=rule, line width=.4pt] (\x,-4) -- (\x,-66);}
    \foreach \y in {-17, -31, -45} {\draw[draw=rule, line width=.4pt] (0,\y) -- (146,\y);}
  \end{scope}
\end{tikzpicture}
\actorlegend
\caption{Vienas veiksmas, dvi pareigos, trys ataskaitos. Kiekvienas stulpelis yra vienas žmogaus
žingsnis; eilutės rodo, ką apie kiekvieną pareigą sako naujausia patvirtinta ataskaita, ką
registras patvirtino ir kokia veiksmo būsena. Kiekvieną šviesiai mėlyną langelį apskaičiuoja
kodas, patvirtinant ataskaitą arba atšaukiant patvirtinimą.}\label{fig:register:timeline}
\end{figure}

Skaitant stulpelį po stulpelio:

\begin{enumerate}
  \item \textbf{Patvirtinama 1~ataskaita.} Jos plane yra veiksmas dviem išvadoms: pareiga X yra
    \status{missing}, pareiga Y --- \status{partial}. Veiksmas įtraukiamas kaip \code{open} su
    dviem elementais.
  \item \textbf{Atsakingas asmuo pažymi jį įvykdytu.} Įrašomas \code{done_at}. Nieko
    nepatvirtinama: tai pareiškimas.
  \item \textbf{Patvirtinama po to parengta 2~ataskaita.} Ji randa X \status{met}, o Y vis dar
    \status{partial}. Elementą dėl X patvirtina 2~ataskaita. Elemento dėl Y --- ne, ir kadangi
    2~ataskaita parengta po \code{done_at}, veiksmas grįžta į \code{in_progress}. Patvirtinimas
    grąžina \code{confirmed: 1}, \code{reopened: 1}, o kai 2~ataskaitos planas vėl nurodo savo
    atvirą išvadą dėl Y, tą išvadą skaičiuoja kaip \code{covered}. Jei 2~ataskaita Y apskritai nebūtų vertinusi, veiksmas
    būtų likęs \code{done}.
  \item \textbf{Patvirtinama 3~ataskaita.} Ji randa abi įvykdytas. Y patvirtina 3~ataskaita; X
    lieka patvirtinta 2~ataskaitos --- pirmosios jos nenutrūkusios sekos ataskaitos. Patvirtintas
    kiekvienas elementas, todėl veiksmą uždaro 3~ataskaita, iš būsenos \code{in_progress}.
  \item \textbf{3~ataskaitos patvirtinimas atšaukiamas} --- ranka arba todėl, kad peržiūrėtojas
    pakeitė vieną iš jos išvadų. Be 3~ataskaitos Y neturi išspręsto vertinimo: jos patvirtinimas
    išvalomas, ir veiksmas grįžta į \code{in_progress}. X lieka patvirtinta, nes 2~ataskaita vis
    dar patvirtinta. Vėl patvirtinus 3~ataskaitą, veiksmas vėl uždaromas.
\end{enumerate}

\subsection{Lemia ataskaitų parengimo, o ne patvirtinimo tvarka}

Ataskaitos ne visada patvirtinamos ta tvarka, kuria buvo parengtos. Kadangi \code{_settle}
vertinimus rikiuoja pagal ataskaitos sukūrimo laiką ir perskaito juos visus, sena, pavėluotai
patvirtinta ataskaita negali nusverti naujesnės. Testas
\code{test_the_reports_speak_in_the_order_they_were_made} ketvirtą ataskaitą, kuri viską randa
įvykdyta, patvirtina anksčiau už trečią, kuri vieną pareigą vis dar randa atvirą: ketvirtoji
veiksmą uždaro, o trečioji, patvirtinta po jos, nieko nepakeičia ir neprideda veiksmo išvadai,
kurią naujesnė ataskaita randa įvykdytą. Tada tas pats testas ataskaitas trina: ištrynus trečią,
uždarymas išlieka, nes ketvirtoji rado tą patį; ištrynus ir ketvirtą, uždarymą pagrindžiantis
įrodymas dingsta, o kartu su juo --- ir uždarymas.

\subsection{Ko atšaukimas neatšaukia}

Veiksmai, kuriuos ataskaita įtraukė, lieka, kai jos patvirtinimas atšaukiamas ar ataskaita
ištrinama: kodo žodžiais, \enquote{dabar jie yra įmonės}. Įrodymą seka tik patvirtinimai ir jais
pagrįsti uždarymai. Veiksmas, kurio ataskaita ištrinta, išsaugo savo elementus ir praranda tik
nuorodą į ataskaitą. Vėl patvirtinus ataskaitą, nieko neįtraukiama du kartus.


\section{Tapatybė: pareiga atpažįstama pagal jos identifikatorių}\label{sec:register:identity}

Viskas, kas aprašyta \cref{sec:register:settle}, sieja vienos ataskaitos išvadą su kitos
ataskaitos išvada per \code{obligationId}. Sąsaja tiksli, be jokio derinimo pagal pavadinimą ar
nuostatą; ji veikia tik tol, kol pareiga iš vienos ataskaitos į kitą išlaiko savo
identifikatorių. Kada taip yra, sprendžiama už registro ribų --- tuo, kaip saugomos pareigos.

\begin{itemize}
  \item \textbf{Teisės akto ar standarto pareiga} išlaiko savo identifikatorių, kai jos aktas
    perskaitomas iš naujo, jei žmogus ją jau patvirtino. Pakartotinis išskyrimas ištrina ir
    pakeičia tik tas eilutes, kurios vis dar \code{proposed} (\file{corpus/obligations.py});
    patvirtintos ir atmestos eilutės niekada neliečiamos. Pasiūlyta pareiga, išskirta dar kartą,
    yra nauja eilutė su nauju identifikatoriumi. Išskyrimas ir peržiūra aprašyti
    \cref{ch:obligations}.
  \item \textbf{Įmonės vidaus reikalavimas} išlaiko savo identifikatorių, kai jo politika
    perskaitoma iš naujo, o jį nustatanti ištrauka nepasikeitė (\code{read_requirements} faile
    \file{analysis/internal.py}): nauja eilutė perima pasiūlytos eilutės su ta pačia cituojama
    ištrauka identifikatorių. Pakeista ištrauka duoda naują reikalavimą.
\end{itemize}

Kai nepatvirtinto elemento pareigos nebėra, jokia vėlesnė ataskaita niekada nebeturės jos
identifikatoriaus, ir veiksmas galėtų laukti patvirtinimo amžinai. Užuot laukęs, registras tai
pasako. \code{standing(actions)} patikrina, kurios elementų pareigos tebeegzistuoja, o
\code{serialize_company_action} nepatvirtintą elementą, kurio pareigos nebėra, pateikia su
\code{obligationId: null}. Tada veiksmas turi \code{byHandOnly: true}, kaip ir kiekvienas
veiksmas, neturintis jokių elementų. Konsolė tokiam veiksmui parodo vieną nekrintančią į akis
eilutę --- \enquote{Šio veiksmo ataskaita patvirtinti negali --- jį uždaro tik žmogus, pastaboje
nurodęs priežastį} --- o prie elemento paaiškina kodėl: nuo to laiko iš naujo perskaitytas jo aktas, o įmonės
vidaus reikalavimo atveju --- jo politika. Saugomas elementas
identifikatorių išlaiko; jį slepia tik atsakymas.

Todėl pareigų registro peržiūra turi būti atlikta prieš pasikliaujant patvirtinimais, o akto,
kurio pasiūlytas pareigas apima atviri veiksmai, nereikėtų išskirti iš naujo, kol jos
neperžiūrėtos (\cref{ch:ops}). 2026~m. lapkričio 13~d. veikiančios sistemos pareigų registre buvo 5\,960
pareigų, iš jų 5\,938 patvirtintos ir 12 vis dar pasiūlytų; kaip ta peržiūra atlikta, nurodyta
\cref{ch:obligations}.


\section{Istorija ir pokyčiai}\label{sec:register:history}

Dvi mažesnės funkcijos naudoja tą pačią sąsają, kad nuo ataskaitos pažvelgtų atgal. Kitaip nei
patvirtinimai, jos skaito kiekvieną ankstesnę \emph{baigtą} įmonės ataskaitą, patvirtintą ar
ne: jos aprašo, ką sakė ataskaitos, o ne ką priėmė žmogus.

\paragraph{Išvados istorija.} \code{history(report)} kiekvienai ataskaitos išvadų pareigai
grąžina paskutinius kartus, kai ją vertino ankstesnė baigta ataskaita, naujausius pirma, ne
daugiau kaip \code{HISTORY}~=~4 įrašus \code{{reportId, date, status}} su galiojančia būsena.
Ataskaita, kuri pareigos nevertino, įrašo nepalieka. \code{serialize_report} sąrašą prideda prie
kiekvienos išsamios ataskaitos išvados kaip \code{history}; jis apskaičiuojamas, kai ataskaita
pateikiama, todėl atspindi vėlesnius peržiūrėtojų sprendimus ankstesnėse ataskaitose. Konsolė
pirmuosius tris įrašus rodo išvados kortelėje kaip ankstesnių vertinimų eilutę, kiekvieną --- su
nuoroda į tą ataskaitą.

\paragraph{Kas pasikeitė.} \code{changes(report)} palygina ataskaitą su ankstesnėmis baigtomis
įmonės ataskaitomis, pareigą po pareigos, ir grąžina tris sąrašus, kiekvieną kaip skaičių su
ne daugiau kaip penkiais pavadinimais. Čia \emph{spraga} yra išvada, kurios būsena
\status{missing} ar \status{partial}; \status{unknown} nėra nei spraga, nei išspręsta.

\begin{table}[H]
\caption{Kaip \code{changes} klasifikuoja ataskaitos pareigą}\label{tab:register:changes}
\begin{tabularx}{\textwidth}{@{}L{20mm}L{44mm}Y@{}}
\toprule
\hd{Raktas} & \hd{Ankstesnis vertinimas → dabar} & \hd{Reiškia} \\
\midrule
\code{closed} & spraga → išspręsta & Spraga, kuri dabar įvykdyta. \\
\code{back} & išspręsta → spraga & Pareiga, kuri buvo įvykdyta, o dabar vėl yra spraga. \\
\code{new} & nei spraga, nei išspręsta arba niekada nevertinta → spraga & Spraga pareigoje,
  kurios nevertino nė viena ankstesnė ataskaita arba kuri anksčiau buvo \status{unknown}. \\
\bottomrule
\end{tabularx}
\end{table}

Ankstesnis vertinimas yra \emph{paskutinės ankstesnės ataskaitos, kuri tą pareigą vertino},
vertinimas, o ne tiesiog prieš tai buvusios ataskaitos. Priežastis praktinė: ataskaita gali būti
parengta keliems dokumentams ar keliems aktams. Jei kiekviena pareiga būtų lyginama su tiesiog
prieš tai buvusia ataskaita, tarp jų parengta siaura ataskaita lemtų, kad kitoje išsamioje ataskaitoje viskas,
ko ji nevertino, atrodytų nauja. Testas
\code{test_a_narrow_report_in_between_does_not_make_the_rest_look_new} šią elgseną įtvirtina.
Atsakymas vis tiek nurodo tiesiog prieš tai buvusią ataskaitą (\code{since}, \code{sinceDate})
kaip datą, nuo kurios kalbama. Išvados dėl įmonės vidaus reikalavimų į pokyčius neįtraukiamos,
nes jos visur skaičiuojamos atskirai, o kai įmonė neturi ankstesnės baigtos ataskaitos, funkcija
nieko negrąžina.

Rezultatas naudojamas trijose vietose, ir visas jas aprašo \cref{ch:report}. Ataskaitai
baigiantis jis išsaugomas ataskaitos statistikoje kaip \code{stats.changes} (\code{report_stats}
faile \file{analysis/pipeline.py}). Jis pridedamas prie santraukos etapo įvesties kaip blokas su
antrašte \texttt{SINCE THE PREVIOUS REPORT}, su trimis skaičiais ir jų pavadinimais, o santraukos
instrukcijai nurodoma pirmąją pastraipą, bendrąjį įvertinimą, užbaigti vienu sakiniu apie
pokyčius. Be to, ataskaitos PDF po skyriumi \enquote{Trumpai} išspausdina vieną eilutę
(\file{analysis/pdf.py}, \code{_glance}): \enquote{Nuo ankstesnės ataskaitos (data): pašalinta
spragų -- n, naujų spragų -- n, vėl atsiradusių -- n.} Visose trijose vietose skaičius
apskaičiuoja kodas; modelis tik suformuluoja sakinį.


\section{Teisės aktų pokyčiai}\label{sec:register:law}

Ataskaita yra įvertinimas pagal tas aktų redakcijas, kurios galiojo jos vykdymo metu. Kai aktas
vėliau pasikeičia, dalis ataskaitos išvadų remiasi nuostatomis, kurios nebeskamba taip, kaip
skambėjo. Įkėlimas užfiksuoja, kas pasikeitė, o registro \code{law_changes} įvardija tas
išvadas.

\paragraph{Kas užfiksuojama.} Kai įkėlimo užduotis aktyvuoja naują akto redakciją, ji lauke
\code{legal_document_versions.changes} išsaugo, ką ta redakcija pakeitė, palyginti su redakcija,
kurią ji pakeičia (\code{version_changes} faile \file{corpus/ingest.py}): \code{{from, changed, added,
removed}}, kur trijuose sąrašuose yra nuostatų keliai akto tvarka. Tekstas, kuriame tarpai normalizuoti,
lyginamas straipsnis po straipsnio --- akto be straipsnių atveju pagal jo aukščiausio lygio
punktus: \code{changed} yra keliai, esantys abiejose redakcijose su skirtingu tekstu,
\code{added} --- esantys tik naujojoje, \code{removed} --- esantys tik senojoje. Pirmoji akto
redakcija įrašo neturi (null); jo neturi ir redakcija, kuri apdorojama dar kartą vietoje, nes ji
būtų lyginama su savimi. Abiejų redakcijų joks modelis neskaito; tai saugomo teksto palyginimas.
Įkėlimas ir redakcijos aprašyti \cref{ch:corpus}.

\paragraph{Ką pateikia maršrutas.} \route{GET}{/api/companies/{id}/law-changes}
(\code{law_changes}) paima naujausią baigtą įmonės ataskaitą ir kiekvienam aktui, kurį ta
ataskaita vertino (kaip \code{assessed} ar \code{upcoming}) ir kurio dabartinė redakcija nebėra ta,
kurią ataskaita skaitė, sujungia visų nuo to laiko atsiradusių redakcijų įrašus. Kiekvienam aktui
grąžinamas ataskaitos skaitytos redakcijos numeris ir dabartinės redakcijos numeris su jos data,
trys sąrašai ir \code{findings}: ataskaitos išvados pagal tą aktą, kurių nuorodos remiasi pakeista
ar panaikinta nuostata, kiekviena su identifikatoriumi, pavadinimu, galiojančia būsena ir cituojamu
keliu. Nuoroda remiasi straipsniu, kai jos kelias yra straipsnio kelias arba jį tęsia be dar vieno
skaitmens, todėl \texttt{Art.\ 68(1)} remiasi \texttt{Art.\ 68}, o \texttt{Art.\ 680} --- ne.
Pridėtos nuostatos jokios esamos išvados neliečia ir pateikiamos informacijai.

Požymis \code{recorded} yra false, kai tarpinė redakcija įrašo neturi. Tada aktas pasikeitė, o
sistema negali pasakyti kur, ir pasako būtent tai, užuot rodžiusi tuščius sąrašus, kurie būtų
perskaityti kaip \enquote{niekas nepasikeitė}: konsolė išspausdina \enquote{Aktas pasikeitė, bet
pakeistos nuostatos neužfiksuotos}, o pokalbio įrankiui pasakoma, kad iš naujo peržiūrėti reikia
visą aktą.

Konsolė sąrašą rodo po registru, dalyje \ui{Teisės aktų pokyčiai}, kiekvieną išvadą susiedama
su jos kortele ataskaitoje. Sąrašas nieko nekeičia: jis neatidaro veiksmo ir nepakeičia išvados.
Jis įvardija, kur žmogui verta pažiūrėti, o nauja analizė įmonę įvertina pagal dabar galiojantį
tekstą.

2026~m. lapkričio 13~d. jokia veikiančios sistemos saugoma redakcija neturėjo pokyčių įrašo: nuo
tada, kai pradėta fiksuoti, nebuvo aktyvuota nė viena nauja akto redakcija. Mechanizmą išbando
testas \code{test_a_new_version_of_an_act_names_the_findings_to_look_at_again}, kuris
perrašo vieną straipsnį, vieną pašalina ir vieną prideda, o tada tikrina tris sąrašus ir
įvardytas išvadas; su tikru pakeitimu mechanizmas dar nesusidūrė.


\section{Registro skaitymas: pokalbis ir Excel eksportas}\label{sec:register:reading}

\subsection{Pokalbių asistentas skaito ir niekada nerašo}

Pokalbių asistentas registrui turi vieną įrankį, \code{getActionRegister}, asistento paslaugos
faile \file{llm/tools/actions.js}. Jis kviečia du skaitymo maršrutus --- registrą, tada teisės
aktų pokyčius --- ir grąžina glaustą atsakymą: skaičius pagal rodomą būseną; ne daugiau kaip
\code{MAX_ACTIONS}~=~20 veiksmų, uždarytus --- tik tada, kai prašoma būtent tos būsenos,
kiekvieną su būsena, atsakingu asmeniu, terminu ir ne daugiau kaip \code{MAX_ITEMS}~=~6
išvadomis, kurias jis uždaro; \code{closedOnlyByHand} ten, kur API sako \code{byHandOnly}; ir,
jei aktai pasikeitė, teisės aktų pokyčius su išvadomis, kurias reikia patikrinti iš naujo. Jei
užklausa į teisės aktų pokyčių maršrutą nepavyksta, ta dalis praleidžiama, o registras vis tiek grąžinamas.

Įrankis nepriima jokio argumento, kuris galėtų ką nors pakeisti, ir joks kitas pokalbio įrankis
nepatvirtina ataskaitos ir nekeičia veiksmo; testas faile \file{tests/tools-actions.test.mjs}
tikrina, kad vienintelis įrankis, kurio pavadinime minimas veiksmas, patvirtinimas ar registras,
yra šis. Tai, ką modelis gali \emph{pasakyti}, apribota dvejopai: įrankio aprašu ir jo
užuomina \code{next}, grąžinama su kiekvienu atsakymu, ir pokalbio instrukcijos taisykle.

\begin{promptexcerpt}{chat.assistant}{13 taisyklė, apie registrą}
[…]
The register moves only when a person signs a report off on the report page: its plan comes in, and a later signed-off report that finds an action's obligations met closes it. You cannot open, close or edit an action, and you never call one closed because a document was uploaded or the user says it is done – "done" there means carried out and awaiting that re-test; tell the user to sign the next report off, or to close the action by hand with a reason on the company's Actions tab.
[…]
\end{promptexcerpt}

Ši taisyklė reikalinga todėl, kad viliojantis atsakymas yra klaidingas. Naudotoją, ką tik įkėlusį
trūkstamą procedūrą ir klausiantį \enquote{ar dabar jau uždaryta?}, \enquote{taip} pradžiugintų, ir
modelis, linkęs padėti, galėtų taip atsakyti. Registro atsakymas yra toks: įkeltas dokumentas
nėra įrodymas, kol ataskaita jo neįvertino ir žmogus tos ataskaitos nepatvirtino. Pokalbio
įrankiai ir instrukcija kaip visuma aprašyti \cref{ch:chat}.

Vieną netiesioginį kelią verta įvardyti. Pokalbis gali įrašyti peržiūrėtojo sprendimą dėl
išvados naudotojo nurodymu ir su naudotojo nurodyta priežastimi (įrankis \code{overrideFinding},
\cref{ch:chat}). Jis siunčia tą pačią užklausą kaip konsolės peržiūros dialogas, todėl
patvirtintai ataskaitai pasekmė ta pati: patvirtinimas atšaukiamas, o tai, ką jis patvirtino
registre, atšaukiama kartu. Patvirtinti ataskaitos pokalbis vis tiek negali.

\subsection{Excel eksportas}

\route{GET}{/api/companies/{id}/actions/export} grąžina visą registrą kaip darbaknygę
(\code{render_actions_xlsx} faile \file{analysis/exports.py}), lietuvių ar anglų kalba pagal
parametrą \code{language}, išdėstytą kurso išvadų ir veiksmų registro stulpeliais. Viena eilutė
yra vienas veiksmas, sąrašo tvarka: eilės numeris; veiksmas su aprašu; išvados, kurias jis turi
uždaryti, kiekviena atskiroje eilutėje su aktu ir nuostata ir žodžiu \enquote{patvirtinta}, kai
ją patvirtino vėlesnė ataskaita; atsakingas asmuo; terminas; rodoma būsena, įskaitant
vėlavimą; prioritetas; ataskaitos, iš kurios veiksmas kilo, data; uždarymo data; kaip uždaryta
(\enquote{patvirtino vėlesnė ataskaita} arba \enquote{uždarė asmuo}); ir pastaba. Kiekvienas
tekstinis langelis saugomas kaip tekstas, niekada kaip formulė. Atsisiuntimas įrašomas į veiksmų
žurnalą kaip \code{actions.export}. Pokalbio įrankis tai, ką veiksmas uždaro, formuluoja taip pat
kaip darbaknygė, todėl abu sutampa.


\section{Kur nubrėžta riba}\label{sec:register:line}

\begin{codebox}{Registre modelis niekur nekviečiamas}
\file{analysis/register.py} importuoja ataskaitų ir teisės aktų bazės modelius, duomenų bazės
transakciją ir laikrodį. Jis neimportuoja funkcijos, kuri kviečia modelį, ir nė viena jo kviečiama
funkcija jos neimportuoja. Plano įtraukimas, išvados patvirtinimas, veiksmo uždarymas, grąžinimas į darbą,
\code{late} apskaičiavimas, išvados istorija, pokyčiai tarp ataskaitų ir teisės aktų pokyčiai ---
visa tai yra saugomų identifikatorių, būsenų, datų ir tekstų palyginimai. Turint tas pačias
ataskaitas, tuos pačius parašus ir tuos pačius žmonių pakeitimus ta pačia tvarka, registras
visada išeina toks pat.

\smallskip
Kas kita yra jo \emph{įvestys}, ir jas reikia įvardyti. Veiksmo pavadinimus, atsakingo asmens
rolę, laikotarpį ir kategoriją ataskaitos plane parašė modelis, o išvados būsena yra modelio
vertinimas, nebent ją pakeitė peržiūrėtojas. Registras joms neprideda jokio savo sprendimo: jis
tik užfiksuoja, kurias iš jų žmogus patvirtino.
\end{codebox}

\begin{humanbox}{Ką čia daro tik žmonės}
Žmogus patvirtina ataskaitą ir gali patvirtinimą atšaukti. Žmogus yra atsakingas už veiksmą,
nustato jo terminą ir nurodo, kokia jo padėtis. Žmogus gali uždaryti veiksmą ranka, nurodydamas
priežastį, kuri lieka įraše, ir gali vėl atidaryti tai, ką uždarė žmogus. Žmogus ginčija
ataskaitos atliktą uždarymą ties išvada, kuria tas uždarymas remiasi. Registre niekas nevyksta
sistemos iniciatyva: kiekvienas pakeitimas išplaukia iš vieno iš šių žmogaus žingsnių.
\end{humanbox}

Riba nubrėžta čia dėl trijų priežasčių. Pirmoji --- atskaitomybė: veiksmų registras pasako, ką
įmonė įsipareigojo padaryti ir ar įsipareigojimą įvykdė, ir abu šie teiginiai turėtų vesti iki
įvardyto žmogaus ir datuotos ataskaitos, o ne iki modelio formuluotės vieną dieną. Antroji ---
atkuriamumas: modelis, dukart paklaustas, ar veiksmas uždarytas, gali dukart atsakyti skirtingai,
o sąsaja pagal pareigos identifikatorių --- ne. Trečioji --- sprendimas, kuriam modelis
\emph{tinka}, jau priimtas, patikrintas ir peržiūrėtas kitur --- išvadoje (\cref{ch:report}).
Registras tą sprendimą panaudoja pakartotinai ir savo neprideda, todėl neprideda ir naujos
vietos, kur modelis galėtų suklysti. Sistemos kontrolės priemones surenka \cref{ch:control} ir
tarp jų aptaria šią.


\section{Kas išbandyta, kas testuojama ir žinomas apribojimas}\label{sec:register:tried}

\begin{lessonbox}{Išbandyta duomenų bazės kopijoje 2026 m. lapkričio 12 d.}
Registras buvo paleistas su tikromis pavyzdinės įmonės \demo{} ataskaitomis duomenų bazės
kopijoje, be jokio modelio iškvietimo. Patvirtinus išsamią įmonės ataskaitą, buvo įtraukti 25 jos
plano veiksmai su 587 atviromis išvadomis, kurias jie turi uždaryti, ir kiekviena turėjo savo
pareigos identifikatorių. Patvirtinus vėlesnę nedidelę ataskaitą, buvo pridėti 7 veiksmai, o 6 praleisti kaip jau vykdomi. Patvirtinus atvirkštine tvarka --- pirma naujesnę ataskaitą --- registre buvo
tos pačios 653 išvados, ir nė vienos neapėmė du veiksmai.

\smallskip
Tada buvo surežisuota pakartotinė patikra: išsami ataskaita nukopijuota, o dalis jos išvadų
pažymėta kaip įvykdytos. Patvirtinus kopiją, buvo patvirtintos 7 išvados, uždarytas veiksmas,
kurio visos išvados buvo įvykdytos, ir į darbą grąžintas vienas veiksmas, kuris buvo pažymėtas
įvykdytu, bet vis dar turėjo atvirą išvadą. Atšaukus tą patvirtinimą, būtent tai ir buvo
atšaukta. Po antros pakartotinės patikros ištrynus kopiją, uždarymai išliko, remdamiesi antrąja.
Tai vienos dienos stebėjimai su viena įmone; pakartotinė patikra buvo surežisuota redaguojant
išvadas, o ne gauta iš antros tikros analizės po tikro trūkumų šalinimo.
\end{lessonbox}

\paragraph{Testai.} \file{tests/test_register.py} turi septynis testus, kurie registrą išbando
per API maršrutus su modelių imitacijomis: patvirtinimas planą įtraukia vieną kartą; vėlesnė
patvirtinta ataskaita yra pakartotinė patikra; ataskaitos kalba ta tvarka, kuria buvo parengtos;
žmogus yra atsakingas už veiksmą, nustato terminą, perkelia ir uždaro veiksmą, ir eksportas;
išvados istorija ir ataskaitos pokyčiai; tarpinė siaura ataskaita; ir nauja akto redakcija.
Prieigos teisių matrica (\file{tests/test_authz_matrix.py}) apima registro maršrutus, o asistento
paslauga testuoja pokalbio įrankį. Testų rinkiniai aprašyti \cref{ch:testing}.

\paragraph{Žinomas apribojimas.} \code{standing()} klausia tik, ar pareigos eilutė tebeegzistuoja.
Tačiau ataskaita vertina dabartinės akto redakcijos pareigas, kurios nėra atmestos. Todėl pareiga,
kurią administratorius atmeta jau po to, kai ją apėmė veiksmas, tebeegzistuoja, niekada
nebevertinama ir nepažymima \code{byHandOnly}: jos veiksmas laukia patvirtinimo, kuris negali
ateiti, kol žmogus tai pastebi ir uždaro jį ranka. 2026~m. lapkričio 13~d. nė vienos iš dešimties
atmestų pareigų neapėmė joks veiksmas. Tas pats, tuo pačiu kodo skaitymu, galioja ir patvirtintai
pareigai, kurią nauja jos akto redakcija paliko nuošalyje, nes pasikeitė jos cituojama nuostata:
\code{carry_forward} tokią eilutę palieka senojoje redakcijoje, kur jos neskaito jokia vėlesnė
ataskaita. Antruoju atveju, kol įmonei nebus parengta kita ataskaita, aktas turi naujesnę redakciją, nei skaitė
jos naujausia ataskaita, todėl pokytis bent jau pateikiamas tarp teisės aktų pokyčių
(\cref{sec:register:law}), kur žmogus į jį nukreipiamas.
